\subsection{为什么 open source 之后，配方会越来越像工程规范}

开源模型的一个重要后果，是数据配方不再只属于单一实验室。任何人都可以比较：你是更多依赖网页、更多依赖百科，还是更多依赖代码和论文；你是保守过滤还是激进过滤；你是追求规模还是追求可复现。这个比较一旦成立，整个领域就会开始建立共同语言，比如 token 数、域占比、去重策略、语言识别和 PII 清理。

这也是为什么很多公开模型的说明文件越来越像系统设计文档，而不是简单的数据来源表。数据工程一旦变成公共话题，大家讨论的就不只是”用了什么”，还包括”为什么选这个””怎么证明这个选择有效”。

